% p026_b 的电脑重画（三角形导体的等效电阻：电流竖直向上/向下 1Ω；两块并联 R并=0.5Ω；沿斜边方向 R=1Ω；水平方向 R未知；上下串联 R串=2Ω）
\begin{tikzpicture}[line width=0.7pt, >={Stealth[length=1.7mm]}]
  % (1) 三角形，电流向上
  \draw (0,0) -- (1.05,0) -- (0.525,0.66) -- cycle;
  \draw[->] (0.3,-0.3) -- (0.3,0.37);
  \draw[->] (0.525,-0.3) -- (0.525,0.58);
  \draw[->] (0.75,-0.3) -- (0.75,0.37);
  \node at (0.525,-0.62) {$1\,\Omega$};
  % (2) 同一三角形，电流向下
  \draw (1.5,0) -- (2.55,0) -- (2.025,0.66) -- cycle;
  \draw[->] (1.8,0.37) -- (1.8,-0.3);
  \draw[->] (2.025,0.58) -- (2.025,-0.3);
  \draw[->] (2.25,0.37) -- (2.25,-0.3);
  \node at (2.025,-0.62) {$1\,\Omega$};
  % (3) 两个三角形拼成平行四边形，电流向上（并联）
  \draw (3.0,0) -- (4.05,0) -- (4.575,0.66) -- (3.525,0.66) -- cycle;
  \draw (3.525,0.66) -- (4.05,0);
  \draw[->] (3.25,-0.3) -- (3.25,0.31);
  \draw[->] (3.65,-0.3) -- (3.65,0.66);
  \draw[->] (4.0,-0.3) -- (4.0,0.66);
  \node at (3.3,-0.62) {$1\,\Omega$};
  \node[anchor=west] at (4.55,0.33) {$1\,\Omega$};
  \node[anchor=west] at (4.1,-0.62) {$R_{\text{并}}=0.5\,\Omega$};
  % (4) 三角形，电流沿斜边方向
  \draw (6.2,0) -- (7.25,0) -- (6.725,0.66) -- cycle;
  \draw[->] (6.245,0.339) -- (6.627,-0.141);
  \draw[->] (6.403,0.537) -- (6.942,-0.141);
  \draw[->] (6.56,0.735) -- (7.257,-0.141);
  \node[anchor=north west] at (7.05,-0.38) {$R=1\,\Omega$};
  % (5) 三角形，电流水平
  \draw (8.55,0) -- (9.6,0) -- (9.075,0.66) -- cycle;
  \draw[->] (8.4,0.165) -- (9.75,0.165);
  \draw[->] (8.4,0.33) -- (9.75,0.33);
  \draw[->] (8.4,0.495) -- (9.75,0.495);
  \node at (9.075,-0.5) {$R$未知};
  % (6) 两个三角形尖对尖上下叠放，电流向上（串联）
  \draw (10.35,0) -- (11.35,0) -- (10.85,0.6) -- cycle;
  \draw (10.35,1.2) -- (11.35,1.2) -- (10.85,0.6) -- cycle;
  \draw[->] (10.65,-0.4) -- (10.65,1.45);
  \draw[->] (11.05,-0.4) -- (11.05,1.45);
  \node[anchor=west] at (11.5,0.95) {$1\,\Omega$};
  \node[anchor=west] at (11.5,0.4) {$1\,\Omega$};
  \node[anchor=west] at (12.3,0.3) {$R_{\text{串}}=2\,\Omega$};
\end{tikzpicture}
